\begin{longtable}{|p{4cm}|p{11cm}|}
    \hline
    \textbf{ID - Nome dello UC} & 1 - Creazione di una issue \cite{uc1-figma} \\
    \hline
    \textbf{Data di creazione} & 11/10/2025 \\
    \hline
    \textbf{Attore primario} & Utente \\
    \hline
    \textbf{Trigger} & L'utente indica di voler creare una nuova issue \\
    \hline
    \textbf{Descrizione} & L'utente ha trovato un problema e intende aprire una nuova issue per tenere traccia di esso. \\
    \hline
    \textbf{Precondizioni} & \textbf{PRE-1.} L'utente è stato registrato da un amministratore ed ha effettuato correttamente l'autenticazione \newline
    \textbf{PRE-2.} Il sistema ha caricato la pagina di riepilogo delle issues \\
    \hline
    \textbf{Postcondizioni} & \textbf{POST-1.} La nuova issue è salvata nel sistema \newline
    \textbf{POST-2.} La nuova issue è visibile nella vista di riepilogo delle issue \\
    \hline
    \textbf{Scenario di successo} & 1. L'utente indica di voler segnalare una nuova issue (M1) \newline
    2. Il sistema apre la modale di creazione di una nuova issue (M2) \newline
    3. L'utente inserisce il titolo della issue (M3) \newline
    4. L'utente indica i dettagli della issue nella descrizione (M4) \newline
    4a. [Opzionale] L'utente sceglie una priorità per la issue \newline
    4b. [Opzionale] L'utente può allegare un'immagine \newline
    4c. [Opzionale] L'utente sceglie un tipo per la issue \newline
    5. L'utente segnala di aver terminato l'inserimento dei dati richiesti (M5) \newline
    6. Il sistema valida i dati inseriti dall'utente \newline
    7. Il sistema salva la nuova issue \newline
    8. Il sistema mostra un messaggio di successo (M6) \\
    \hline

    \multicolumn{2}{r}{\small Continua nella pagina successiva} \\
    \textbf{Estensioni} & \textbf{4a. Scelta della priorità per la issue:} \newline
    4a1. L'utente clicca sul menù “priorità” (M9) \newline
    4a2. Il sistema apre un menù con le possibili priorità (M10) \newline
    4a3. L'utente sceglie una tra le alternative (M11) \newline
    4a4. Il sistema mostra l'alternativa scelta (M12) \newline
    4a5. Il flusso ritorna al termine del punto 4 dello scenario di successo \newline

    \textbf{4b. Caricamento di un'immagine:} \newline
    4b1. L'utente clicca sul bottone per caricare un'immagine (M13) \newline
    4b2. Il sistema apre il file explorer (M14) \newline
    4b3. L'utente seleziona un'immagine dal proprio dispositivo (M15, M16) \newline
    4b4. Il sistema valida il formato e la dimensione dell'immagine \newline
    4b5. Il sistema sostituisce l'icona di “allega” con un messaggio di successo (\checkmark) e aggiunge un bottone di eliminazione accanto alla label “immagine”(M17) \newline
    4b7. Il flusso torna al termine del punto 4a dello scenario di successo \newline

    \textbf{4b3a. L'utente annulla il caricamento dell'immagine:} \newline
    4b3a1. L'utente clicca sul bottone “annulla” \newline
    4b3a2. Il flusso ritorna al termine del punto 4 dello scenario di successo \newline

    \textbf{4b4a. Errore di validazione dell'immagine:} \newline
    4b4a1. Il sistema mostra un messaggio di errore e invita l'utente a riprovare (M18) \newline
    4b4a2. Il flusso ritorna al termine del punto 4 dello scenario di successo \newline

    \textbf{4b6a. Eliminazione dell'immagine allegata:} \newline
    4b6a1. L'utente clicca sul bottone di eliminazione \newline
    4b6a2. Il sistema elimina l'immagine allegata \newline
    4b6a3. Il sistema sostituisce l'icona di successo (\checkmark) con il bottone “allega” e rimuove il bottone di eliminazione \newline
    4b6a4. Il flusso ritorna al termine del punto 4 dello scenario di successo \\
    \hline
    \multicolumn{2}{r}{\small Continua nella pagina successiva}
    \end{longtable}


    \begin{longtable}{|p{4cm}|p{11cm}|}
    \hline
    \textbf{Estensioni} & 
    \textbf{4c. Scelta del tipo di issue:} \newline
    4c1. L'utente clicca sul menù “tipo” (M19) \newline
    4c2. Il sistema apre il menù con i vari tipi di issues (M20) \newline
    4c3. L'utente sceglie un tipo di issue (M21) \newline
    4c4. Il sistema mostra il tipo scelto (M22) \newline
    4c5. Il flusso ritorna al termine del punto 4 dello scenario di successo \newline

    \textbf{6a. Errore nella validazione dell'input (titolo):} \newline
    6a1. L'utente inserisce un titolo non valido \newline
    6a2. Il sistema evidenzia il campo del titolo e mostra un messaggio all'utente, segnalando la non validità del titolo scelto (M7) \newline
    6a3. L'utente modifica il titolo, inserendone uno corretto (M3) \newline
    6a4. Il flusso ritorna al punto 4 dello scenario di successo \newline

    \textbf{6b. Errore nella validazione dell'input (descrizione):} \newline
    6b1. L'utente inserisce una descrizione non valida \newline
    6b2. Il sistema evidenzia il campo della descrizione e mostra un messaggio all'utente, segnalando la non validità della descrizione scelta (M8) \newline
    6b3. L'utente inserisce una descrizione valida (M4) \newline
    6b4. Il flusso ritorna al punto 5 nello scenario di successo \newline

    \textbf{7a. Errore nel salvataggio della issue:} \newline
    7a1. Il sistema riscontra un errore durante il salvataggio della issue \newline
    7a2. Il sistema mostra uno specifico messaggio di errore, invitando l'utente a riprovare (M23) \newline
    7a3. Il flusso ritorna al termine del punto 4 nello scenario di successo \\
    \hline
    \textbf{Priorità} & Alta \\
    \hline
    \textbf{Frequenza d'uso} & Presumibilmente tutti gli utenti, tra 1 e 5 volte al giorno per utente \\
    \hline
    \textbf{Informazioni associate} & Mostrato nella tabella ~\ref{tab:uc1-proprieta} \\
    \hline
\end{longtable}

\newpage